% !TEX root = ../../ctfp-print.tex

\lettrine[lhang=0.17]{如}{今谈起函数式编程}, 你不可能不提单子. 但在另一个平行宇宙里, Eugenio Moggi
碰巧把注意力投向了 Lawvere 理论, 而不是单子. 让我们去探索一下那个宇宙.

\section{泛代数}

描述代数的方式有很多, 抽象层次各不相同. 我们试图找到一种通用的语言来描述幺半群、群、环这类东西. 在
最简单的层面上, 这些构造都定义了一个集合的元素上的\emph{运算}, 再加上这些运算必须满足的
一些\emph{定律}. 比如, 幺半群可以用一个满足结合律的二元运算来定义. 我们还有单位元和单位律.
但只要稍微发挥一点想象力, 我们就能把单位元变成一个零元运算 —— 一个不取任何参数、只返回集合中
某个特殊元素的运算. 如果我们想谈群, 就再加一个一元运算符, 它取一个元素并返回它的逆. 与之配套的还有
相应的左逆律和右逆律. 一个环定义了两个二元运算符, 外加更多定律. 如此等等.

大图景是这样的: 一个代数由一族 $n$ 元运算 (取各个不同的 $n$ 值) 和一族等式恒等式来定义.
这些恒等式全都是全称量化的. 结合律方程必须对三个元素的所有可能组合都成立, 以此类推.

顺便一提, 这把域排除在外了, 理由很简单: 零 (关于加法的单位元) 关于乘法没有逆. 域所要求的逆律无
法全称量化.

只要把运算 (函数) 换成态射, 这个泛代数的定义就能推广到 $\Set$ 之外的范畴. 我们不再选取一个集合, 而
是选取一个物件 $a$ (叫作泛物件). 一元运算就只是 $a$ 的一个自态射. 那别的元数呢
(\newterm{arity 元数}指的是某个运算所取参数的个数)? 二元运算 (元数为 2) 可以定义成一个从
积$a\times{}a$ 回到 $a$ 的态射. 一般的 $n$ 元运算是一个从 $a$ 的 $n^\text{th}$ 次幂到 $a$ 的态射:
\[\alpha_n \Colon a^n \to a\] 零元运算是一个从终止物件出发的态射 (也就是 $a$ 的零次幂). 所以,
要定义任何代数, 我们只需要一个范畴, 它的物件是某个特殊物件 $a$ 的各次幂. 具体的代数则编码在这个范畴的
hom-集合里. 这就是一言以蔽之的 Lawvere 理论.

Lawvere 理论的推导要经过许多步骤, 所以这里给出一张路线图:

\begin{enumerate}
  \tightlist
  \item
有限集合构成的范畴 $\cat{FinSet}$.
  \item
它的骨架 $\cat{F}$.
  \item
它的反范畴 $\Fop$.
  \item
Lawvere 理论 $\cat{L}$: 范畴 $\cat{Law}$ 里的一个物件.
  \item
Lawvere 范畴的模型 $M$: 范畴\\ $\cat{Mod}(\cat{Law}, \Set)$ 里的一个物件.
\end{enumerate}

\begin{figure}[H]
  \centering
  \includegraphics[width=0.8\textwidth]{images/lawvere1.png}
\end{figure}

\section{Lawvere 理论}

所有 Lawvere 理论都共享一个共同的骨架. 一个 Lawvere 理论里的所有物件, 都是由一个物件经积 (其实就是
幂)生出来的. 但在一般的范畴里, 我们怎么定义这些积? 结果发现, 我们可以借助一个从更简单的范畴出发的
映射来定义积. 事实上, 这个更简单的范畴可以定义的是余积而不是积, 而我们再用一个\emph{反变}函子把
它们嵌入目标范畴. 反变函子把余积变成积, 把入射变成投影.

Lawvere 范畴的骨架, 最自然的选择就是有限集合构成的范畴 $\cat{FinSet}$. 它包含空集
$\varnothing$、单元素集合 $1$、两元素集合 $2$, 等等. 这个范畴里的所有物件都可以由
单元素集合经余积生成(把空集当成零元余积的特例). 比如, 两元素集合是两个单元素集合的和, $2 = 1 + 1$,
用Haskell 表达就是:

\src{snippet01} 然而, 尽管我们很自然地认为空集只有一个, 单元素集合却可以有许多个互不相同的. 特别地,
集合$1 + \varnothing$ 不同于集合 $\varnothing + 1$, 也不同于 $1$ —— 尽管它们全都同构. 集合范畴里的
余积并不满足结合律. 我们可以通过构造一个把同构的集合都认同为一的范畴来补救. 这样的
范畴叫作\newterm{skeleton 骨架}. 换句话说, 任何 Lawvere 理论的骨架就是 $\cat{FinSet}$ 的骨架$\cat{F}$
.
这个范畴里的物件可以与自然数 (包括零) 相认同, 这些自然数对应于 $\cat{FinSet}$ 里元素的个数.
余积扮演加法的角色. $\cat{F}$ 里的态射对应于有限集合之间的函数. 比如, 从 $\varnothing$ 到 $n$ 有
唯一的态射 (空集是初始物件), 从 $n$ 到 $\varnothing$ 没有态射 ($\varnothing \to \varnothing$ 除外),
从 $1$ 到 $n$ 有 $n$ 个态射 (那些入射), 从 $n$ 到 $1$ 有一个态射, 如此等等. 这里 $n$ 表示$\cat{F}$
里的一个物件, 它对应于 $\cat{FinSet}$ 里所有 $n$ 元素集合经同构认同之后的东西.

借助范畴 $\cat{F}$, 我们可以正式地把 \newterm{Lawvere theory} 定义成一个装备了特殊函子的范畴$\cat{L}$
:
\[I_{\cat{L}} \Colon \Fop \to \cat{L}\] 这个函子必须在物件上是双射, 并且必须保持有限积 ($\Fop$ 里的
积等同于 $\cat{F}$ 里的余积): \[I_{\cat{L}}\ (m\times{}n) = I_{\cat{L}}\ m\times{}I_{\cat{L}}\ n\]
你有时会看到有人把这个函子刻画成在物件上是恒等的, 意思就是 $\cat{F}$ 与 $\cat{L}$ 里的物件是同一批.
因此我们对它们使用同样的名字 —— 就用自然数来记. 但要记住, $\cat{F}$ 里的物件并不等同于集合 (它们是
同构集合的类).

一般来说, $\cat{L}$ 里的 hom-集合比 $\Fop$ 里的更丰富. 它们可能含有不对应于 $\cat{FinSet}$ 里
任何函数的态射 (后者有时叫作 \newterm{basic product operations 基本积运算}). Lawvere 理论的
等式定律就编码在这些态射里.

关键的观察是: $\cat{F}$ 里的单元素集合 $1$ 被映到某个物件上, 我们在 $\cat{L}$ 里也叫它 $1$; 而
$\cat{L}$ 里所有别的物件都自动是这个物件的幂. 比如, $\cat{F}$ 里的两元素集合 $2$ 是余积 $1 + 1$,
所以它必定被映到 $\cat{L}$ 里的积 $1 \times 1$ (或者 $1^2$). 在这个意义上, 范畴 $\cat{F}$ 的行为就
像$\cat{L}$ 的对数.

$\cat{L}$ 的态射当中, 有一批是由函子 $I_{\cat{L}}$ 从 $\cat{F}$ 搬过来的. 它们在 $\cat{L}$ 里
扮演结构性的角色. 特别地, 余积入射 $i_k$ 变成了积投影 $p_k$. 一个有用的直觉, 是把
投影\[p_k \Colon 1^n \to 1\] 想象成一个 $n$ 元函数的原型, 它把除第 $k^\text{th}$ 个变量以外的所有
变量都忽略掉. 反过来, $\cat{F}$里的常值态射 $n \to 1$ 变成了 $\cat{L}$ 里的对角态射 $1 \to 1^n$.
它们对应于变量的复制.

$\cat{L}$ 里有意思的态射, 是那些定义投影之外的 $n$ 元运算的态射. 正是这些态射把一个 Lawvere 理论与
另一个区分开来. 它们是乘法、加法、单位元的选取, 等等, 也就是定义该代数的那些东西. 但要让 $\cat{L}$
成为一个完整的范畴, 我们还需要复合运算 $n \to m$ (等价地, $1^n \to 1^m$). 由于这个范畴的结构很简单,
它们结果都是更简单的 $n \to 1$ 型态射的积. 这是一句话的推广: 返回一个积的函数是若干函数的积 (或者,
如我们先前见过的, Hom-函子是连续的).

\begin{figure}[H]
  \centering
  \includegraphics[width=0.8\textwidth]{images/lawvere1.png}
    \caption{Lawvere 理论 $\cat{L}$ 建立在 $\Fop$ 之上, 并从它那里继承那些定义积的
    ``无聊的'' 态射. 它再加上描述 $n$ 元运算的 ``有意思的'' 态射 (虚线箭头).}
\end{figure}

Lawvere 理论构成一个范畴 $\cat{Law}$, 其中的态射是保持有限积、并且与那些函子 $I$ 可交换的函子.
给定两个这样的理论 $(\cat{L}, I_{\cat{L}})$ 与 $(\cat{L'}, I'_{\cat{L'}})$, 它们之间的一个态射是
一个函子 $F \Colon \cat{L} \to \cat{L'}$, 使得:
\begin{gather*}
  F\ (m \times n) = F\ m \times F\ n \\
  F \circ I_{\cat{L}} = I'_{\cat{L'}}
\end{gather*}
Lawvere 理论之间的态射, 把将一个理论解释到另一个理论内部这一想法封装了起来. 比如, 若我们忽略掉逆,
群的乘法就可以被解释成幺半群的乘法.

最简单的平凡例子是 Lawvere 范畴 $\Fop$ 本身 (对应于把恒等函子选作 $I_{\cat{L}}$). 这个既没有运算也
没有定律的 Lawvere 理论, 恰好是 $\cat{Law}$ 里的初始物件.

讲到这里, 如果能看到一个非平凡的 Lawvere 理论例子会很有帮助, 但在弄清楚模型是什么之前, 很难把
它解释清楚.

\section{Lawvere 理论的模型}

理解 Lawvere 理论的关键在于意识到: 一个这样的理论泛化了大量具有相同结构的个体代数. 比如, 幺半群的
Lawvere 理论描述了成为一个幺半群这件事的本质. 它必须对所有幺半群都成立. 某个具体的幺半群就成为
这样一个理论的模型. 模型被定义成一个从 Lawvere 理论 $\cat{L}$ 到集合范畴 $\Set$ 的函子. (还有
一些推广的 Lawvere 理论用别的范畴来当模型, 但这里我只关注 $\Set$.) 由于 $\cat{L}$ 的结构严重依赖积,
我们要求这样的函子保持有限积. 于是 $\cat{L}$ 的一个模型 —— 也叫作 Lawvere 理论 $\cat{L}$ 上的代数
——由这样一个函子来定义:
\begin{gather*}
  M \Colon \cat{L} \to \Set \\
  M\ (a \times b) \cong M\ a \times M\ b
\end{gather*}
注意, 我们只要求积在\emph{同构的意义下}被保持. 这一点非常重要, 因为严格地保持积会把大多数有意思的
理论都排除掉.

模型保持积这件事意味着, $M$ 在 $\Set$ 里的像是这样一串集合: 它们由集合 $M\ 1$ 的各次幂生成 —— 而
$M\ 1$ 是物件 $1$ 在 $\cat{L}$ 里的像. 我们把这个集合叫作 $a$. (这个集合有时被叫作\emph{类} (sort),
而这样的代数叫作 \newterm{single-sorted 单类的}. 还有一些把 Lawvere 理论推广到多类代数的做法.) 特别地
,
$\cat{L}$ 里的二元运算被映成函数: \[a \times a \to a\] 和任何函子一样, $\cat{L}$ 里的多个态射有
可能被压成 $\Set$ 里的同一个函数.

顺便说一句, 所有定律都是全称量化的等式, 这一事实意味着每个 Lawvere 理论都有一个平凡模型: 一个常函子,
它把所有物件都映到单元素集合, 把所有态射都映到它上面的恒等函数.

$\cat{L}$ 里一个形如 $m \to n$ 的一般态射被映成一个函数: \[a^m \to a^n\] 如果我们有两个不同的模型
$M$与 $N$, 它们之间的一个自然变换是一族以 $n$ 为下标的函数: \[\mu_n \Colon M\ n \to N\ n\] 或者等价地
:
\[\mu_n \Colon a^n \to b^n\] 其中 $b = N\ 1$.

注意, 自然性条件保证了 $n$ 元运算被保持: \[N\ f \circ \mu_n = \mu_1 \circ M\ f\] 其中
$f \Colon n \to 1$ 是 $\cat{L}$ 里的一个 $n$ 元运算.

定义模型的那些函子构成一个模型范畴 $\cat{Mod}(\cat{L}, \Set)$, 其中的态射是自然变换.

考虑平凡 Lawvere 范畴 $\Fop$ 的一个模型. 这样的模型完全由它在 $1$ 处的取值 $M\ 1$ 决定. 由于
$M\ 1$可以是任何集合, 这种模型要多少有多少, 与 $\Set$ 里集合的数量一样多. 而且,
$\cat{Mod}(\Fop, \Set)$ 里的每个态射 (也就是函子 $M$ 与 $N$ 之间的自然变换) 都由它在 $M\ 1$ 处的
分量唯一确定. 反过来, 每个函数 $M\ 1 \to N\ 1$ 都诱导出两个模型 $M$ 与 $N$ 之间的一个自然变换.
因此$\cat{Mod}(\Fop, \Set)$ 等价于 $\Set$.

\section{幺半群理论}

Lawvere 理论最简单的非平凡例子, 描述的是幺半群的结构. 它是单独一门理论, 却提炼出了所有可能的
幺半群的结构, 其意义在于: 这门理论的模型铺满了整个幺半群范畴 $\cat{Mon}$. 我们先前
已经见过一个\hyperref[free-monoids]{普遍构造}, 它表明每个幺半群都能从某个合适的自由
幺半群出发、通过把一族态射认同为一而得到. 所以单独一个自由幺半群就已经泛化了大量的幺半群. 不过,
自由幺半群有无穷多个. 幺半群的 Lawvere 理论 $\cat{L}_{\cat{Mon}}$ 把它们全都收进了一个优雅的构造里.

每个幺半群都必须有单位元, 所以 $\cat{L}_{\cat{Mon}}$ 里必须有一个从 $0$ 到 $1$ 的特殊态射 $\eta$.
注意 $\cat{F}$ 里不可能有对应的态射: 这样的态射方向相反, 从 $1$ 到 $0$; 而在 $\cat{FinSet}$ 里,
那就是一个从单元素集合到空集的函数. 这样的函数不存在.

接着考虑态射 $2 \to 1$, 也就是 $\cat{L}_{\cat{Mon}}(2, 1)$ 的成员, 它们之中必须含有所有二元运算的原型
.
在 $\cat{Mod}(\cat{L}_{\cat{Mon}}, \Set)$ 里构造模型时, 这些态射会被映成从笛卡尔积
$M\ 1 \times M\ 1$到 $M\ 1$ 的函数. 换句话说, 就是有两个自变量的函数.

问题在于: 只用幺半群运算符, 我们能实现多少个二元函数? 我们把这两个自变量叫作 $a$ 与 $b$. 有
一个函数忽略两个自变量、直接返回幺半群的单位元. 然后有两个投影, 分别返回 $a$ 和 $b$. 再接下去是
返回$ab$、$ba$、 $aa$、$bb$、$aab$ 等等\ldots{} 的各个函数. 事实上, 这种二元函数的数量, 等于以 $a$ 与
$b$ 为生成元的自由幺半群里元素的个数. 注意 $\cat{L}_{\cat{Mon}}(2, 1)$ 必须含有所有这些态射, 因为
其中一个模型就是自由幺半群. 在自由幺半群里, 它们对应着互不相同的函数. 别的模型可能把
$\cat{L}_{\cat{Mon}}(2, 1)$ 里的多个态射压成同一个函数, 但自由幺半群不会.

如果我们把以 $n$ 个元素为生成元的自由幺半群记作 $n^*$, 那么可以把 hom-集合 $\cat{L}(2, 1)$ 与
幺半群范畴 $\cat{Mon}$ 里的 hom-集合 $\cat{Mon}(1^*, 2^*)$ 相认同. 一般地,
我们取$\cat{L}_{\cat{Mon}}(m, n)$ 为 $\cat{Mon}(n^*, m^*)$. 换句话说, 范畴 $\cat{L}_{\cat{Mon}}$ 就
是自由幺半群范畴的反范畴.

幺半群 Lawvere 理论的\emph{模型}范畴\\ $\cat{Mod}(\cat{L}_{\cat{Mon}}, \Set)$, 等价于所有
幺半群构成的范畴 $\cat{Mon}$.

\section{Lawvere 理论与单子}

你可能还记得, 代数理论可以用单子来描述 —— 特别是 \hyperref[algebras-for-monads]{单子的代数}.
那么Lawvere 理论与单子之间存在联系, 就毫不奇怪了.

首先我们看看, 一个 Lawvere 理论如何诱导出一个单子. 它是通过一个
\hyperref[free-forgetful-adjunctions]{伴随}来实现的, 这个伴随位于一个遗忘函子和一个自由函子之间.
遗忘函子 $U$ 给每个模型指派一个集合. 这个集合的来历, 是把 $\cat{Mod}(\cat{L}, \Set)$ 里的函子 $M$
求值到 $\cat{L}$ 里的物件 $1$ 上.

推导 $U$ 还有另一种方式, 靠的是 $\Fop$ 是 $\cat{Law}$ 里初始物件这一事实. 这意味着, 对任何 Lawvere
理论 $\cat{L}$, 都存在唯一一个函子 $\Fop \to \cat{L}$. 这个函子在模型上诱导出反向的函子 (因为
模型是\emph{从}理论到集合的函子): \[\cat{Mod}(\cat{L}, \Set) \to \cat{Mod}(\Fop, \Set)\]
但正如我们讨论过的, $\Fop$ 的模型范畴等价于 $\Set$, 于是我们就得到了那个遗忘函子:
\[U \Colon \cat{Mod}(\cat{L}, \Set) \to \Set\] 可以证明, 这样定义出来的 $U$ 总有左伴随, 也就是自由
函子 $F$.

对有限集合来说这很容易看出来. 自由函子 $F$ 生成自由代数. 一个自由代数, 是
$\cat{Mod}(\cat{L}, \Set)$里由一个有限的生成元集合 $n$ 生成的特殊模型. 我们可以把 $F$ 实现成可表函子:
\[\cat{L}(n, -) \Colon \cat{L} \to \Set\] 要说明它确实是自由的, 我们只需证明它是遗忘函子的左伴随:
\[\cat{Mod}(\cat{L}(n, -), M) \cong \Set(n, U(M))\] 我们把右边化简一下:
\[\Set(n, U(M)) \cong \Set(n, M\ 1) \cong (M\ 1)^n \cong M\ n\] (我用到了态射集合同构于指数这一事实,
而在这里那个指数就是迭代出来的积.) 这个伴随是米田引理的结果:
\[[\cat{L}, \Set](\cat{L}(n, -), M) \cong M\ n\] 遗忘函子与自由函子合在一起, 就在 $\Set$ 上定义了
一个\hyperref[monads-categorically]{单子} $T = U \circ F$. 于是每个 Lawvere 理论都生成一个单子.

结果发现, \hyperref[algebras-for-monads]{这个单子的代数}范畴等价于模型的范畴.

你可能还记得, 单子代数定义了求值那些用单子构造出来的表达式的方法. Lawvere 理论定义了
可用来生成表达式的 $n$ 元运算. 而模型则提供了对这些表达式求值的手段.

不过, 单子与 Lawvere 理论之间的联系并不是双向的. 只有有限元单子才能引向 Lawvere 理论. 有
限元单子建立在有限元函子之上. $\Set$ 上的一个有限元函子完全由它在有限集合上的作用决定. 它在
任意集合$a$ 上的作用, 可以用下面这个余端求值: \[F\ a = \int^n a^n \times (F\ n)\] 由于余端泛化了
余积、或者说和, 这个公式是幂级数展开的推广. 或者我们也可以用函子是一个推广的容器这个直觉. 那样的话,
一个装着 $a$ 的有限元容器可以描述成形状与内容的和. 这里 $F\ n$ 是用于存放$n$个元素的形状的集合, 而
内容是一个 $n$ 元组 —— 它本身又是 $a^n$ 的一个元素. 比如, 列表 (作为一个函子) 就是有限元的:
每一种元数对应一个形状. 树的每个元数上有更多形状, 如此等等.

首先, 所有由 Lawvere 理论生成的单子都是有限元的, 它们都可以表达成余端:
\[T_{\cat{L}}\ a = \int^n a^n \times \cat{L}(n, 1)\] 反过来, 给定 $\Set$ 上任何有限元单子 $T$,
我们都能构造出一个 Lawvere 理论. 我们从构造 $T$ 的 Kleisli范畴开始. 你可能还记得, Kleisli 范畴里从
$a$ 到 $b$ 的一个态射, 是由底层范畴里的一个态射给出的 : \[a \to T\ b\] 当限制在有限集合上时, 这就变成
:
\[m \to T\ n\] 与这个 Kleisli 范畴相反的那个范畴 $\cat{Kl}^{op}_{T}$, 限制在有限集合上, 就是
我们要找的那个Lawvere理论. 特别地, 描述 $\cat{L}$ 里 $n$ 元运算的 hom-集合 $\cat{L}(n, 1)$, 由
hom-集合$\cat{Kl}_{T}(1, n)$ 给出.

结果发现, 我们在编程里遇到的大多数单子都是有限元的, 一个显眼的例外是延续单子. 我们有可能把 Lawvere
理论的概念推广到有限元运算之外.

\section{作为余端的单子}

我们把那个余端公式再仔细看一遍. \[T_{\cat{L}}\ a = \int^n a^n \times \cat{L}(n, 1)\] 首先,
这个余端是在 $\cat{F}$ 上的一个前函子 $P$ 上取的, 它定义为: \[P\ n\ m = a^n \times \cat{L}(m, 1)\]
这个前函子对第一个参数 $n$ 是反变的. 我们来看看它是怎么提升态射的. $\cat{FinSet}$ 里的一个态射是有
限集合之间的映射 $f \Colon m \to n$. 这样的映射描述的是从一个 $n$ 元素集合里挑出 $m$ 个元素 (允许重复
).
它可以被提升为 $a$ 的各次幂之间的映射, 也就是 (注意方向): \[a^n \to a^m\] 这个提升只不过是从 $n$
个元素的元组 $(a_1, a_2,...a_n)$ 里挑出 $m$ 个元素 (可能带重复).

\begin{figure}[H]
  \centering
  \includegraphics[width=0.5\textwidth]{images/liftpower.png}
\end{figure}

\noindent
比如, 取 $f_k \Colon 1 \to n$ —— 从一个 $n$ 元素集合里挑出第 $k^\text{th}$ 个元素. 它提升成一个函数,
取 $a$ 的一个 $n$ 元组并返回其中第 $k^\text{th}$ 个.

或者取 $f \Colon m \to 1$ —— 一个把所有 $m$ 个元素都映到同一个元素上的常值函数. 它的提升是一个函数,
取 $a$ 的单个元素并把它复制 $m$ 次: \[\lambda{}x \to (\underbrace{x, x,\ ...\ , x}_{m})\]
你可能会注意到, 要说清这个前函子对第二个参数是协变的, 并不是一眼就能看出来的.
hom-函子$\cat{L}(m, 1)$实际上对 $m$ 是反变的. 不过我们取余端的范畴不是 $\cat{L}$ 而是 $\cat{F}$.
余端变量 $n$ 遍历的是有限集合 (或者它们的骨架). 范畴 $\cat{L}$ 含有 $\cat{F}$ 的反范畴, 所以
$\cat{F}$ 里的一个态射 $m \to n$ 是 $\cat{L}$ 里 $\cat{L}(n, m)$ 的成员 (这个嵌入由函子
$I_{\cat{L}}$给出).

我们来检验一下 $\cat{L}(m, 1)$ 作为从 $\cat{F}$ 到 $\Set$ 的函子的函子性.
我们想提升一个函数$f \Colon m \to n$, 所以目标是从 $\cat{L}(m, 1)$ 到 $\cat{L}(n, 1)$ 实现一个函数.
对应于函数 $f$ 的是 $\cat{L}$ 里一个从 $n$ 到 $m$ 的态射 (注意方向). 把这个态射与 $\cat{L}(m, 1)$
做前复合, 就得到了 $\cat{L}(n, 1)$ 的一个子集.

\begin{figure}[H]
  \centering
  \begin{tikzcd}[column sep=large]
    \cat{L}(m, 1) \arrow[r] & \cat{L}(n, 1)\\
    {}^m \bullet \arrow[r, "f"'] & \bullet^n
  \end{tikzcd}
\end{figure}

\noindent
注意, 通过提升一个函数 $1 \to n$, 我们能从 $\cat{L}(1, 1)$ 走到 $\cat{L}(n, 1)$. 稍后我们会用到这一点
.

一个反变函子 $a^n$ 与一个协变函子 $\cat{L}(m, 1)$ 的积, 是一个前函子 $\Fop \times \cat{F} \to \Set$.
记住, 余端可以定义成前函子所有对角成员的不交并 (不交和), 其中某些元素被认同为一. 这些认同对
应着余楔条件.

在这里, 余端起先是所有 $n$ 上集合 $a^n \times \cat{L}(n, 1)$ 的不交并. 那些认同可以通过把
\hyperref[ends-and-coends]{余端表达成余等化子}来生成. 我们从一个非对角项
$a^n \times \cat{L}(m, 1)$开始. 要走到对角上, 我们可以把一个态射 $f \Colon m \to n$ 作用到积的
第一个或第二个分量上. 得到的两个结果随即被认同为一.

\begin{figure}[H]
  \centering
  \begin{tikzcd}
    & a^n \times \cat{L}(m, 1)
    \arrow[dl, "\langle f {,} \id \rangle"']
    \arrow[dr, "\langle \id {,} f \rangle"]
    & \\
    a^m \times \cat{L}(m, 1)
    & \scalebox{2.5}[1]{\sim}
    & a^n \times \cat{L}(n, 1) \\
    & f \Colon m \to n &
  \end{tikzcd}
\end{figure}

\noindent
我先前已经展示过, 提升 $f \Colon 1 \to n$ 会得到这两个变换: \[a^n \to a\] 以及:
\[\cat{L}(1, 1) \to \cat{L}(n, 1)\] 因此, 从 $a^n \times \cat{L}(1, 1)$ 出发, 我们两条路都能走到:
\[a \times \cat{L}(1, 1)\] 这是提升 $\langle f, \id \rangle$ 的时候; 而
当提升$\langle \id, f \rangle$ 时, 我们得到: \[a^n \times \cat{L}(n, 1)\] 但这并不
意味着$a^n \times \cat{L}(n, 1)$ 的所有元素都能与 $a \times \cat{L}(1, 1)$ 认同为一. 那是因为
$\cat{L}(n, 1)$ 里并非所有元素都能从 $\cat{L}(1, 1)$ 到达. 记住, 我们只能提升来自 $\cat{F}$ 的态射.
$\cat{L}$ 里一个非平凡的 $n$ 元运算, 是没法通过提升一个态射 $f \Colon 1 \to n$ 构造出来的.

换句话说, 在余端公式里, 我们只能认同那些 $\cat{L}(n, 1)$ 可以通过施加基本态射从 $\cat{L}(1, 1)$ 到
达的求和项. 它们全都等价于 $a \times \cat{L}(1, 1)$. 基本态射就是那些来自 $\cat{F}$ 的态射的像.

我们看看最简单的情形 —— Lawvere 理论 $\Fop$ 本身 —— 会怎样. 在这样的理论里, 每个 $\cat{L}(n, 1)$ 都
能从 $\cat{L}(1, 1)$ 到达. 这是因为 $\cat{L}(1, 1)$ 是一个只含恒等态射的单元素集合, 而
$\cat{L}(n, 1)$ 只含有对应于 $\cat{F}$ 里入射 $1 \to n$ 的那些态射, 它们\emph{确实}是基本态射.
因此余积里的所有求和项都彼此等价, 于是我们得到: \[T\ a = a \times \cat{L}(1, 1) = a\] 而这就是
恒等单子.

\section{副作用的 Lawvere 理论}

既然单子与 Lawvere 理论之间有如此紧密的联系, 那么很自然地要问: 在编程里, Lawvere 理论能不能被
用作单子的替代品? 单子的主要问题是它们不能很好地复合. 我们没有一个通用的配方来打造单子变换器.
Lawvere理论在这方面有优势: 它们可以用余积和张量积来复合. 反过来说, 只有有
限元单子才能轻松地转成Lawvere 理论. 这里的异类是延续单子. 这个方向上还有研究在进行 (见参考文献).

为了让你尝尝用 Lawvere 理论描述副作用是什么味道, 我来讨论一个简单的例子: 传统上用 \code{Maybe}
单子实现的异常.

\code{Maybe} 单子是由带一个零元运算 $0 \to 1$ 的 Lawvere 理论生成的. 这门理论的一个模型, 是一个把
$1$映到某个集合 $a$、并把那个零元运算映成一个函数的函子:

\src{snippet02} 我们可以用余端公式把 \code{Maybe} 单子找回来. 我们来看看加入这个零元运算对
hom-集合$\cat{L}(n, 1)$做了什么. 除了造出一个新的 $\cat{L}(0, 1)$ ($\Fop$ 里没有它) 之外,
它还往$\cat{L}(n, 1)$ 里添加了新的态射. 这些是形如 $n \to 0$ 的态射与我们的 $0 \to 1$ 复合的结果.
在余端公式里, 这类贡献全都被认同为 $a^0 \times \cat{L}(0, 1)$, 因为它们可以这样得到:
\[a^n \times \cat{L}(0, 1)\] 办法是用两种不同的方式提升 $0 \to n$.

\begin{figure}[H]
  \centering
  \begin{tikzcd}
    & a^n \times \cat{L}(0, 1)
    \arrow[dl, "\langle f {,} \id \rangle"']
    \arrow[dr, "\langle \id {,} f \rangle"]
    & \\
    a^0 \times \cat{L}(0, 1)
    & \scalebox{2.5}[1]{\sim}
    & a^n \times \cat{L}(n, 1) \\
    & f \Colon 0 \to n &
  \end{tikzcd}
\end{figure}

\noindent
余端就化简成了: \[T_{\cat{L}}\ a = a^0 + a^1\] 或者, 用 Haskell 记法:

\src{snippet03} 它等价于:

\src{snippet04} 注意, 这门 Lawvere 理论只支持抛出异常, 不支持处理异常.

\section{挑战}

\begin{enumerate}
  \tightlist
  \item
枚举 $\cat{F}$ (即 $\cat{FinSet}$ 的骨架) 里 $2$ 与 $3$ 之间的所有态射.
  \item
证明: 幺半群 Lawvere 理论的模型范畴, 等价于列表单子的单子代数范畴.
  \item
幺半群的 Lawvere 理论生成列表单子. 证明它的二元运算可以用相应的 Kleisli 箭头生成.
  \item
\textbf{FinSet} 是 $\Set$ 的子范畴, 而且有一个函子把它嵌入 $\Set$. $\Set$ 上的任何函子都可以限制到
$\cat{FinSet}$ 上. 证明: 一个有限元函子等于它自身之限制的左 Kan 扩张.
\end{enumerate}

\section{延伸阅读}
\begin{enumerate}
  \tightlist
  \item
\urlref{http://www.tac.mta.ca/tac/reprints/articles/5/tr5.pdf}{Functorial Semantics of Algebraic Theories},
F. William Lawvere
  \item
\urlref{http://homepages.inf.ed.ac.uk/gdp/publications/Comp_Eff_Monads.pdf}{Notions of computation determine monads},
Gordon Plotkin and John Power
\end{enumerate}
